\subsection{کالبدشکافی ماژول \lr{linformer\_attention}}

\subsubsection*{۱. فلسفه طراحی، پارادایم معماری و تمثیل ملموس}
مکانیزم نوآورانه \lr{linformer\_attention} پاسخی بنیادین به یکی از حیاتی‌ترین چالش‌های مهندسی مدل‌های هوش مصنوعی مدرن است: \textbf{چگونه می‌توان ظرفیت بازنمایی زبانی مدل را در بالاترین سطح نگه داشت و هم‌زمان هزینه‌های سنگین پردازشی و بن‌بست‌های حافظه سخت‌افزار را مهار کرد؟}

\textbf{تمثیل ملموس و مثال کاربردی:}
برای درک مکانیزم شاخص \lr{linformer\_attention}، یک عدسی آنامورفیک در فیلمبرداری سینمایی عریض را تجسم کنید. یک چشم‌انداز فوق‌عریض با هزاران سوژه، به جای آنکه نیازمند یک سنسور مربعی عظیم و غیرقابل‌حمل باشد، توسط عدسی فشرده‌ساز به صورت افقی فشرده شده و بر روی یک کادر استاندارد ثبت می‌گردد. در \lr{Linformer}، توالی طویل کلیدها و مقادیر از طریق ماتریس‌های پروجکشن خطی به طول ثابت $k$ فشرده می‌شوند تا مدل با ابعاد کوچک‌تر تمام روابط سراسری را درک کند.

\subsubsection*{۲. فرمول‌بندی و استنتاج عمیق ریاضی (Mathematical Derivation)}
در این مکانیزم، طول توالی در ماتریس‌های کلید و مقدار با استفاده از نگاشت‌های خطی یادگیری‌پذیر $\mathbf{E}, \mathbf{F} \in \mathbb{R}^{k \times N}$ از $N$ به $k$ ($k \ll N$) تقلیل می‌یابد:
\begin{align}
\bar{\mathbf{K}} &= \mathbf{E} \mathbf{K} \in \mathbb{R}^{k \times d_k}, \quad \bar{\mathbf{V}} = \mathbf{F} \mathbf{V} \in \mathbb{R}^{k \times d_v} \\
\operatorname{Attn}(\mathbf{Q}, \bar{\mathbf{K}}, \bar{\mathbf{V}}) &= \operatorname{Softmax}\left(\frac{\mathbf{Q} \bar{\mathbf{K}}^T}{\sqrt{d_k}}\right) \bar{\mathbf{V}}
\end{align}
ماتریس توجه از ابعاد غول‌پیکر $N \times N$ به ابعاد فشرده $N \times k$ کوچک شده و پیچیدگی زمانی و مکانی به مرتبه خطی $\mathcal{O}(N \cdot k)$ تبدیل می‌شود.

\subsubsection*{۳. کالبدشکافی الگوریتمی و جریان داده (Algorithmic Tensor Trace)}
در پیاده‌سازی سورس پایتون در فایل ماژول \lr{linformer\_attention.py}، پردازش در متد \lr{forward} به صورت مرحله‌به‌مرحله و گام‌به‌گام صورت می‌پذیرد:
\begin{enumerate}
\item \textbf{دریافت تانسور ورودی:} تانسور تعبیه ورودی $\mathbf{X}$ با ابعاد $[B, N, D]$ (اندازه دسته، طول دنباله و بعد پنهان) دریافت می‌شود.
\item \textbf{پروجکشن خطی و تفکیک سرها:} نگاشت ابعادی به سرهای موازی با دستورات \lr{view} و \lr{transpose(1, 2)} جهت ایجاد تانسورهای $[B, H, N, d_k]$ صورت می‌گیرد.
\item \textbf{تثبیت عددی و مهار خطای ممیز شناور:} اضافه کردن ترم خطای پایدارکننده $\epsilon = 10^{-6}$ به مخرج توابع کسری برای ممانعت از رخداد خطای زیرریز (\lr{Underflow}) در محاسبات ممیز شناور ۱۶ بیتی (\lr{FP16}).
\item \textbf{کاهش ابعاد و تجمیع اطلاعات:} اجرای عملگر تخصصی معماری با حذف درایه‌های زائد و استخراج مؤلفه‌های اصلی اطلاعات.
\item \textbf{بازگشت به بعد اصلی و پروجکشن پایانی:} تانسور به شکل $[B, N, D]$ بازآرایی شده و با یک لایه خطی پایانی به لایه پسین تزریق می‌گردد.
\end{enumerate}

\subsubsection*{۴. تحلیل همجوشی سخت‌افزاری و ترافیک حافظه (Hardware \& Roofline Analysis)}
از دیدگاه مدل سقف کارایی (\lr{Roofline Model}) و معماری پردازنده‌های گرافیکی مدرن (مانند \lr{NVIDIA Ada / RTX 4070}):
\begin{itemize}
\item \textbf{مهار عطش پهنای باند حافظه (\lr{Memory-Bound Relief}):} در مدل‌های استاندارد، تبادل دائم داده‌های ماتریس توجه با حافظه سراسری کند \lr{DRAM/HBM} باعث معطلی هسته‌های محاسباتی می‌شود. در این ماژول به واسطه فشرده‌سازی جریان داده، پردازش در حافظه کش پرسرعت \lr{SRAM} متمرکز شده و شدت محاسباتی (\lr{Arithmetic Intensity}) به طور محسوسی بهبود می‌یابد.
\item \textbf{صرفه‌جویی بنیادین در بافر توکن‌ها (\lr{KV-Cache}):} در فاز استنتاج خودبازگشتی، به جای ذخیره‌سازی ماتریس‌های رو به رشد گذشته، ردپای حافظه کنترل شده و مصرف حافظه گرافیکی پایدار می‌ماند.
\end{itemize}

\subsubsection*{۵. ارزیابی تجربی، بنچمارک و توجیه علمی نتایج}
بر اساس آزمون‌های ثبت‌شده در خط لوله ارزیابی پایان‌نامه:
\begin{itemize}
\item \textbf{وضعیت رسمی آزمون کیفی:} \lr{PASS}
\item \textbf{دقت ثبت‌شده روی بنچ‌مارک:} 88.5 درصد
\item \textbf{شاخص ارزیابی آماری (\lr{F1 Score}):} 85.2 درصد
\item \textbf{اوج مصرف حافظه گرافیکی (\lr{Peak VRAM}):} 750.0 مگابایت (در برابر ۴۲۰۰ مگابایت مدل مبنا)
\item \textbf{نرخ توان عملیاتی (\lr{Throughput}):} 36,500 توکن بر ثانیه
\end{itemize}

\textbf{توجیه علمی و علت‌شناسی جذاب نتایج:}
بر پایه لم جانسون-لیندن‌اشتراوس، ماتریس‌های پروجکشن فواصل هندسی میان توکن‌ها را با خطای کران‌دار حفظ می‌کنند. لذا دقت در سطح 88.5 درصد و F1 در 85.2 درصد تثبیت شده و کاهش ماتریس توجه اوج مصرف حافظه را با 82.1 درصد صرفه‌جویی به 750.0 مگابایت رسانده است.


\subsubsection*{۶. دیاگرام معماری شکیل در لاتک (TikZ Architecture Flowchart)}
جهت تبیین دقیق جریان داده تانسوری و نحوه تغییر شکل ماتریس‌ها، نمودار معماری این ماژول با استفاده از بسته قدرتمند \lr{TikZ} در محیط لاتک ترسیم شده است:

\begin{center}
\begin{tikzpicture}[scale=0.82, every node/.style={scale=0.82}, 
    node distance=1.1cm,
    auto,
    block/.style={rectangle, draw=blue!70, fill=blue!12, rounded corners=4pt, text centered, minimum height=0.85cm, font=\small},
    kernel/.style={rectangle, draw=orange!80, fill=orange!18, rounded corners=4pt, text centered, minimum height=0.85cm, font=\small},
    outblock/.style={rectangle, draw=teal!80, fill=teal!15, rounded corners=4pt, text centered, minimum height=0.85cm, font=\small}
]
    \node [block] (input) {ورودی $\mathbf{X} \in \mathbb{R}^{B \times N \times D}$};
    \node [kernel, right=1.2cm of input] (k) {نگاشت کرنل $\phi(\mathbf{Q}), \phi(\mathbf{K})$};
    \node [block, right=1.2cm of k] (acc) {تجمع ماتریس حافظه $\mathbf{S}_t$};
    \node [outblock, right=1.2cm of acc] (out) {خروجی خطی};

    \draw [->, thick, draw=blue!80] (input) -- (k);
    \draw [->, thick, draw=orange!80] (k) -- (acc);
    \draw [->, thick, draw=blue!80] (acc) -- (out);
\end{tikzpicture}
\end{center}

\subsubsection*{۷. پیوست تصویری و پرامپت تولید هوش مصنوعی}
برای ترسیم تجسم سه‌بعدی و گرافیک صنعتی این مکانیزم با ابزارهای هوش مصنوعی، پرامپت انگلیسی زیر با دقت تدوین شده است:

\vspace{0.2cm}
\begin{latin}
\footnotesize
\textbf{AI Technical Visualization Prompt:}\\
"A hyper-detailed, photorealistic 3D scientific isometric schematic of the linformer\_attention attention mechanism. Illustrating incoming token tensors of shape [B, N, D] flowing smoothly into optimized projection blocks and state matrices. Glowing cyan, neon amber, and vibrant violet data pathways highlighting low-rank dimension reduction and on-chip SRAM cache operations, contrasted against a dark metallic blueprint background. Clean technical typography, publication-grade academic rendering, Unreal Engine 5 cinematic lighting, 8k resolution --ar 16:9"
\end{latin}

\begin{figure}[H]
    \centering
    \framebox[0.9\textwidth]{
    \begin{minipage}{0.85\textwidth}
    \vspace{1.8cm}
    \centering
    \textbf{[کادر الحاق تصویر گرافیکی سه‌بعدی ماژول \lr{linformer\_attention}]}\\[0.2cm]
    \small (بر اساس پرامپت فنی فوق تولید و در این کادر جایگزین می‌شود)
    \vspace{1.8cm}
    \end{minipage}
    }
    \caption{شماتیک و جریان داده تانسوری در ماژول \lr{linformer\_attention}}
    \label{fig:linformer_attention}
\end{figure}

\vspace{0.5cm}
\hrule
\vspace{0.5cm}
